\section{Problem formulation}
\label{sec:problem}

\subsection{Setting and notation}
\label{sec:problem-setting}

A bank has $K$ risk types. For risk type $k=1,\dots,K$ the loss over the horizon $h$ is a random variable $L_k$ with distribution
function $F_k$, and the model of the risk type delivers $F_k$ as a quantile function $Q_k(p)=\inf\{x:F_k(x)\ge p\}$. The study
uses $K=4$ (credit, market, interest rate risk in the banking book and operational risk), $h$ of one year and a confidence level
$\alpha=0.999$. The value at risk and the expected shortfall of a loss $X$ at level $\alpha$ are
\begin{equation}
\VaR_\alpha(X)=Q_X(\alpha),\qquad
\ES_\alpha(X)=\frac{1}{1-\alpha}\int_\alpha^1 Q_X(p)\dd p ,
\label{eq:var-es}
\end{equation}
and $\mathrm{EL}_k=\E[L_k]$ is the expected loss of risk type $k$. Table~\ref{tab:notation} lists the main symbols of the paper with their
unit. Money is in billions of Chilean pesos (CLP billion) unless a table states CLP million, the unit of the committed documents.

\begin{table}[htbp]
\centering
\small
\caption{The notation of the formulation and of the proposed method: the main symbols with their meaning and their unit; money is in CLP billion unless a table states CLP million.}
\label{tab:notation}
\fitwidth{\begin{tabular}{lp{9cm}l}
\toprule
Symbol & Meaning & Unit \\
\midrule
$K$, $k$ & \relax number of risks (four) and a risk index & \relax 1 \\
$L_k$, $F_k$ & \relax loss of risk $k$ over the horizon and its distribution function & \relax CLP billion \\
$\alpha$, $h$ & \relax confidence level (0.999) and horizon (one year) & \relax 1, year \\
$S$ & \relax total loss $\sum_k L_k$ & \relax CLP billion \\
$\mathrm{EL}_k$, $\mathrm{VaR}_k$, $\mathrm{ES}_k$ & \relax expected loss, value at risk and expected shortfall of $L_k$ at level $\alpha$ & \relax CLP billion \\
$\mathrm{EC}_k$ & \relax economic capital $\mathrm{VaR}_k-\mathrm{EL}_k$ & \relax CLP billion \\
$R$, $\rho_{ij}$ & \relax correlation matrix of the copulas and its entries & \relax 1 \\
$\nu$ & \relax degrees of freedom of the $t$ copula & \relax 1 \\
$b_j$ & \relax capital of row $j$ of the band & \relax CLP billion \\
$H$ & \relax headline capital, the simple sum $\sum_k\mathrm{EC}_k$ & \relax CLP billion \\
$\mathrm{DF}_j$ & \relax diversification factor $b_j/H$ of row $j$ & \relax 1 \\
$\underline{b}$, $\overline{b}$ & \relax lower bound and worst case of the band & \relax CLP billion \\
$\lambda_\nu(\rho)$ & \relax upper tail dependence coefficient of the bivariate $t$ copula & \relax 1 \\
$\tau$ & \relax Kendall rank correlation & \relax 1 \\
\bottomrule
\end{tabular}}
\end{table}


The four marginals that the study uses are those of Banco Andino, the synthetic bank of the case; they are read from the committed
documents and summarised in Table~\ref{tab:marginals}. Their sizes are very different: the credit and the market marginal carry \bdEcShareCredit{} and \bdEcShareMarket{} percent of the sum of the capital figures, the interest rate marginal \bdEcShareIrrbb{} percent and the operational marginal \bdEcShareOperational{} percent.

\begin{table}[htbp]
\centering
\small
\caption{The four marginals of Banco Andino at one year and 99.9 percent, in CLP billion: expected loss, value at risk, expected shortfall, economic capital (value at risk less expected loss) and the Monte Carlo standard error of the value at risk where the marginal is simulated, read from the committed baseline document.}
\label{tab:marginals}
\fitwidth{\begin{tabular}{lrrrrr}
\toprule
Risk & EL & VaR & ES & EC & VaR s.e. \\
\midrule
\relax credit & \relax 37.6 & \relax 298.2 & \relax 346.9 & \relax 260.6 & \relax 1.53 \\
\relax market & \relax 0.0 & \relax 283.4 & \relax 355.8 & \relax 283.4 & \relax -- \\
\relax IRRBB & \relax 27.9 & \relax 100.7 & \relax 106.4 & \relax 72.8 & \relax 0.28 \\
\relax operational & \relax 1.2 & \relax 42.4 & \relax 171.1 & \relax 41.2 & \relax -- \\
\relax total (simple sum of capital) & \relax 66.7 & \relax 724.8 & \relax 980.2 & \relax 658.0 & \relax -- \\
\bottomrule
\end{tabular}}
\end{table}


\subsection{The capital basis and the headline}
\label{sec:problem-basis}

The capital of a risk type is its economic capital, the unexpected loss at level $\alpha$,
\begin{equation}
\mathrm{EC}_k=\VaR_\alpha(L_k)-\mathrm{EL}_k ,\qquad H=\sum_{k=1}^K \mathrm{EC}_k ,
\label{eq:ec-headline}
\end{equation}
and $H$, the simple sum, is the headline capital of the aggregate. The choice of the simple sum as the headline is a supervisory
position of the case, not a result of this paper: a regulator that does not recognise the diversification between risk types adds
the capital needed for each. Every other figure the protocol produces is information for the validator and is not capital. Two conventions of the basis are
stated here. The expected loss is subtracted because the capital figures of the committed documents are the unexpected part of the
loss, the part that provisions and pricing are not meant to absorb, and the expected loss is additive over risk types, which is
what lets one number, $\sum_k\mathrm{EL}_k$, convert the value at risk of any aggregate into its capital. The value at risk, and
not the expected shortfall, defines the capital, because that is how the documents define it; the expected shortfall enters only
the allocation, where its coherence is needed.

The aggregate that a bank reports below $H$ asserts that the four risk types do not reach their $\alpha$-quantiles together. That
assertion is a statement about the joint law of $(L_1,\dots,L_K)$, which no bank observes in the region of probability $1-\alpha$.

\subsection{Dependence uncertainty}
\label{sec:problem-uncertainty}

Let $\mathcal C(F_1,\dots,F_K)$ be the set of joint laws of $(L_1,\dots,L_K)$ with the marginals $F_k$, which the literature calls the
couplings of the marginals. For a coupling $\pi$ write $S_\pi=\sum_k L_k$ for the total loss under $\pi$. The expected loss of
the total does not depend on $\pi$, because the expectation is additive, so the capital of the aggregate under $\pi$ is
\begin{equation}
b(\pi)=\VaR_\alpha(S_\pi)-\sum_{k=1}^K \mathrm{EL}_k ,\qquad \pi\in\mathcal C(F_1,\dots,F_K).
\label{eq:coupling-capital}
\end{equation}
The set $\{b(\pi)\}$ is the dependence uncertainty of the aggregate: it is what the capital can be when only the marginals are
known. Its extreme values are
\begin{equation}
\underline b=\inf_{\pi\in\mathcal C} b(\pi),\qquad \overline b=\sup_{\pi\in\mathcal C} b(\pi),
\label{eq:sharp-bounds}
\end{equation}
the sharp bounds, and the diversification factor of any row $j$ that reports a capital $b_j$ is $\mathrm{DF}_j=b_j/H$.

The problem of the paper is the following. Given the marginals, a correlation matrix $R$ for the copulas and a set $\mathcal N$ of
degrees of freedom of the $t$ copula, define a reporting rule for the aggregate such that (i) the capital is never shown as one
figure but as a fixed ordered band of rows computed on the same marginals, one horizon and one level; (ii) exactly one row is
capital and the other rows say that they are information; (iii) the sharp bounds of Equation~\ref{eq:sharp-bounds} are rows of the
band, so that the band contains the capital under every coupling; (iv) an allocation of the aggregate to the risk types adds up to
the aggregate; and (v) a program can check that each of these holds, and fail when one does not. The paper then measures how much
of the dependence uncertainty the band spans and how it compares with the single figures that are used in its place.

\subsection{Assumptions}
\label{sec:problem-assumptions}

The study rests on four assumptions, each stated with its source. The marginals are given: the protocol does not estimate them,
and every statement below is conditional on them, so a marginal error moves every row of the band alike and the band is not a
remedy for it. The horizon and the level are common to the four marginals; a marginal read at another level or horizon is
converted by the caller before it enters the band, and the experiments in Section~\ref{sec:ablation} measure what an unconverted
marginal does. The correlation matrices are conventions, given entry by entry in Table~\ref{tab:matrices}: the base matrix has the
market-credit correlation of 0.5, the correlation of 0.2 of each of these with operational risk, and, by the case's own statement,
the same values for interest rate risk as for market risk, which the case records as unverified against any source and which the
band exists to expose; the stress matrix raises the market-credit correlation to 0.8, and the crisis matrix of the planted
structures sets the three financial risks at 0.8 with each other and 0.4 with operational risk. The set $\mathcal N$ is $\{4,8,30\}$, a convention of the case that spans heavy tails to nearly Gaussian ones.
Finally, the marginals are continuous and strictly increasing on the region where they are read, which makes the quantile
function a proper inverse and the value at risk of a coupling well defined; the grids of the committed marginals satisfy this
(Section~\ref{sec:design-data}).

\begin{table}[htbp]
\centering
\small
\caption{The three correlation matrices of the study, one entry per pair of risk types (the matrices are symmetric with a unit diagonal): the base matrix of the committed documents, the matrix of the market-credit stress, and the crisis matrix of two planted structures.}
\label{tab:matrices}
\fitwidth{\begin{tabular}{lrrr}
\toprule
Pair of risks & Base & Market-credit 0.8 & Crisis \\
\midrule
\relax credit and market & \relax 0.50 & \relax 0.80 & \relax 0.80 \\
\relax credit and IRRBB & \relax 0.50 & \relax 0.50 & \relax 0.80 \\
\relax credit and operational & \relax 0.20 & \relax 0.20 & \relax 0.40 \\
\relax market and IRRBB & \relax 0.50 & \relax 0.50 & \relax 0.80 \\
\relax market and operational & \relax 0.20 & \relax 0.20 & \relax 0.40 \\
\relax IRRBB and operational & \relax 0.20 & \relax 0.20 & \relax 0.40 \\
\bottomrule
\end{tabular}}
\end{table}

